\section{Zadanie 5 (Kacper Orszulak)}
Udowodnij, że jeśli $\left\{ \Sigma_\alpha : \alpha \in I \right\}$ jest rodziną $\sigma$-algebr na jakimś ustalonym zbiorze $\Omega$ (dla dowolnego zbioru indeksów $I$), to $\bigcap_{\alpha \in I} \Sigma_\alpha$ też jest $\sigma$-algebrą.

\subsection*{Rozwiązanie}
Niech $\Sigma := \bigcap_{\alpha \in I} \Sigma_\alpha$. Po kolei sprawdzamy warunki.
\begin{enumerate}
    \item \textit{$\emptyset, \Omega \in \Sigma$}
        \begin{quote}
        $\forall_{\alpha \in I} \; \emptyset, \Omega \in \Sigma_\alpha \implies \emptyset, \Omega \in \bigcap_{\alpha \in I} \Sigma_\alpha = \Sigma$
        \end{quote}
    \item \textit{Jeśli $E \in \Sigma$ to $\Omega/E \in \Sigma$}
        \begin{quote}
        Niech $E \in \Sigma$. Wtedy $\forall_{\alpha \in I} \; E \in \Sigma_\alpha \iff \Omega / E \in \Sigma_\alpha$.

        Dostajemy zatem $\Omega/E \in \bigcap_{\alpha \in I} \Sigma_\alpha = \Sigma$.
        \end{quote}
    \item \textit{Jeśli $E_1, ..., E_n \in \Sigma$, to $\bigcup_{i=1}^n E_i \in \Sigma$}
        \begin{quote}
        Niech $E_1, ..., E_n \in \Sigma$. Wtedy $\forall_{\alpha \in I} \; E_1, ..., E_n \in \Sigma_\alpha \implies \bigcup_{i=1}^n E_i \in \Sigma_\alpha$.

        Dostajemy zatem $\bigcup_{i=1}^n E_i \in \bigcap_{\alpha \in I} \Sigma_\alpha = \Sigma$.
        \end{quote}
    \item \textit{Jeśli $E_1, E_2, ... \in \Sigma$, to $\bigcup_{i=1}^\infty E_i \in \Sigma$}
        \begin{quote}
        Niech $E_1, E_2, ... \in \Sigma$. Wtedy $\forall_{\alpha \in I} \; E_1, E_2, ... \in \Sigma_\alpha \implies \bigcup_{i=1}^\infty E_i \in \Sigma_\alpha$.

        Dostajemy zatem $\bigcup_{i=1}^\infty E_i \in \bigcap_{\alpha \in I} \Sigma_\alpha = \Sigma$.
        \end{quote}
\end{enumerate}
$\square$
